\section{Results and Evaluation}
\label{sec:results}

This evaluation is qualitative and behavioral: it characterizes what the system
achieves and how its architecture behaves, rather than reporting formal accuracy
or timing benchmarks, which are identified as future work in
Section~\ref{sec:limitations}.

\subsection{Functional Outcomes}
The system meets its central goal: scanned and image-based PDFs that return
nothing to in-page search become fully searchable in the browser. Specifically:
\begin{itemize}
  \item \textbf{Scanned PDFs become searchable.} Pages with no text layer are
        OCR'd and their text injected, so queries match content that was
        previously invisible to search.
  \item \textbf{Mixed pages are handled.} Pages combining native text with text
        baked into images expose both, with duplicates removed so a word is never
        double-counted.
  \item \textbf{Rotated text aligns.} Because rotation is recovered from the
        transform matrix (Eq.~\eqref{eq:angle}) rather than special-cased,
        highlights and selection track rotated pages and rotated runs.
  \item \textbf{Selection and copy work.} Injected text is real DOM text, so
        users can not only find but also select and copy recovered content;
        search-only correction variants are excluded from copy to avoid
        duplication.
  \item \textbf{Fully offline.} OCR engine, language models, and PDF.js are
        bundled, so no document data leaves the machine and the extension works
        without a network connection.
\end{itemize}

\subsection{Architectural Behavior}
The performance-oriented design choices produce the following observable
behavior:

\begin{table}[ht]
  \centering
  \caption{Design mechanisms and the behavior they produce.}
  \label{tab:results}
  \renewcommand{\arraystretch}{1.3}
  \begin{tabularx}{\textwidth}{@{}lX@{}}
    \toprule
    \textbf{Mechanism} & \textbf{Observed effect} \\
    \midrule
    Lazy, viewport-driven rendering & Large documents open immediately;
      only near-viewport pages do work up front. \\
    Idle-time background indexing & Off-screen pages still become searchable,
      scheduled during idle time so scrolling is never blocked. \\
    Web Worker OCR pool & Recognition runs off the main thread on several pages
      at once; the interface stays responsive during OCR. \\
    High-resolution OCR canvas + preprocessing & Small text inside images is
      legible to the recognizer independent of on-screen zoom. \\
    IndexedDB caching (normalized coords) & A document is OCR'd once; re-opening
      or re-zooming it is instant and correct at any scale. \\
    Correction variants & Common digit/letter misreads remain findable by their
      true spelling. \\
    \bottomrule
  \end{tabularx}
\end{table}

\subsection{Component Complexity}
Table~\ref{tab:scope} summarizes the relative implementation complexity of each
component. The coordinate-mapping and DOM-overlay subsystems dominate: they are
where most of the subtle, correctness-critical logic lives.

\begin{table}[ht]
  \centering
  \caption{Relative complexity by component.}
  \label{tab:scope}
  \renewcommand{\arraystretch}{1.25}
  \begin{tabular}{@{}lc@{}}
    \toprule
    \textbf{Component} & \textbf{Complexity} \\
    \midrule
    PDF rendering (PDF.js)            & Medium \\
    OCR pipeline (Tesseract.js)       & Medium \\
    Chrome extension integration      & Medium \\
    Coordinate mapping                & High \\
    DOM text-overlay system           & High \\
    \bottomrule
  \end{tabular}
\end{table}

\subsection{Discussion}
The results validate the project's central thesis: by converting documents into
real, well-positioned DOM text, the system inherits the browser's mature search
behavior almost for free, and the engineering effort concentrates where the real
difficulty lies---geometric alignment and keeping a CPU-heavy workload off the
critical path. The accuracy layer (high-resolution rendering, preprocessing, and
correction variants) measurably widens what is findable in practice, though it
cannot fully overcome poor source scans; this is examined next.
